\section{Results}
\subsection{Production patterns and the baseline comparison}
Electricity use varied widely within each machine, alongside differences in dosed mass and product allocation (Figure~\ref{fig:production}). The scatter of individual records shows why neither an average total consumption nor a single intensity value can fully describe this production setting. Differences between machine-level distributions also reflect differences in what the machines processed. They cannot be read directly as differences in mechanical efficiency.

\fig{02_production_and_energy}{production}{1}{Production and electricity across the three machines. (a--c) Dosed mass and pelleting electricity for all 2,742 positive-energy records, with machine-specific medians annotated. (d--f) Mass distributions, dosing-batch frequencies and product-family mix for all 2,745 source records, including the three zero-energy records. Printed values inside the product bars are record counts. Scatter points represent individual operations, without an implied uniform time spacing.}

In the original comparison on 800 development records, Ridge and histogram gradient boosting gave nearly identical MAE: 51.48 and 51.50 kWh, respectively (Table~\ref{tab:development}). ExtraTrees was close behind, while the multilayer perceptron and the three reference estimators had larger errors. The recent-intensity estimator was less accurate than the training-period intensity reference, despite using completed observations from the validation weeks. Recent information was therefore useful only if incorporated in a suitable way; recency alone did not provide a better baseline.

\begin{table}[!htbp]\centering\small
\caption{Original chronological development comparison on the same 800 records. Configurations were fixed; computational budgets were not identical. The recent-SEC method additionally updates its information during validation.}\label{tab:development}
\begin{tabular}{@{}lrrrr@{}}\toprule
Method & MAE & RMSE & Bias & WAPE (\%)\\\midrule
Ridge & 51.48 & 72.31 & $-3.14$ & 19.88\\
Gradient boosting & 51.50 & 73.46 & $-5.72$ & 19.88\\
ExtraTrees & 52.31 & 75.37 & $-4.89$ & 20.20\\
MLP, three-seed mean & 56.64 & 81.36 & $-1.47$ & 21.87\\
Machine SEC & 59.69 & 83.92 & $-6.59$ & 23.05\\
Recent SEC, last five records & 73.51 & 109.55 & 4.29 & 28.38\\
Machine median & 82.61 & 112.78 & $-14.09$ & 31.89\\\bottomrule
\end{tabular}\par\vspace{4pt}\footnotesize MAE, RMSE and bias are in kWh per record.
\end{table}


Ridge was retained for its accuracy and transparent structure, then fitted on the 1,542 eligible development records. On the subsequent 798-record test, it reduced MAE from 61.29 kWh for the machine-specific intensity reference to 52.45 kWh, an improvement of 8.84 kWh or 14.42\% (Table~\ref{tab:test}). The paired 95\% block-bootstrap interval for the MAE reduction was [6.08, 11.50] kWh. RMSE also decreased, from 92.05 to 80.42 kWh. Removing the five test records with durations exceeding 24 hours left MAE at 52.32 kWh for Ridge and 61.15 kWh for the reference, so those records did not account for the overall advantage.

The improvement occurred on all three machines, but their remaining errors were uneven. Pellet~2 had the lowest Ridge MAE, at 43.36 kWh, compared with 53.36 kWh for Pellet~3 and 66.37 kWh for Pellet~1. Pellet~1 also had a mean prediction error of $-53.85$ kWh, indicating substantial underestimation. The observed-versus-predicted plots make this visible: many of its high-consumption records lie below the equality line (Figure~\ref{fig:testscatter}). Thus, the global improvement coexists with a machine-specific limitation that matters for operational use.

\begin{table}[!htbp]\centering\small
\caption{Archived original test results. Lower MAE for Ridge is observed in each machine; the larger negative bias on Pellet~1 remains visible.}\label{tab:test}
\begin{tabular}{@{}llrrrr@{}}\toprule
Population & Method & $N$ & MAE & RMSE & Bias\\\midrule
All machines & Ridge & 798 & 52.45 & 80.42 & $-21.89$\\
All machines & Machine SEC & 798 & 61.29 & 92.05 & $-24.57$\\\addlinespace
Pellet 1 & Ridge & 200 & 66.37 & 105.14 & $-53.85$\\
Pellet 1 & Machine SEC & 200 & 72.31 & 113.90 & $-64.03$\\
Pellet 2 & Ridge & 333 & 43.36 & 61.26 & $-10.85$\\
Pellet 2 & Machine SEC & 333 & 49.29 & 73.24 & $-8.66$\\
Pellet 3 & Ridge & 265 & 53.36 & 80.10 & $-11.63$\\
Pellet 3 & Machine SEC & 265 & 68.05 & 94.79 & $-14.79$\\\bottomrule
\end{tabular}\par\vspace{4pt}\footnotesize Error metrics are in kWh per record. Global WAPE is 19.11\% for Ridge and 22.34\% for SEC.
\end{table}


\fig{04_original_test_scatter}{testscatter}{1}{Frozen Ridge predictions and errors for every record in the original chronological test. (a--c) Observed versus predicted electricity on identical axes, with equality dashed and MAE, RMSE and bias reported in kWh. (d--f) Signed errors by origin date, with zero dashed and each machine's mean error in red. The 200 Pellet~1, 333 Pellet~2 and 265 Pellet~3 records account for all 798 test observations. The residual axes preserve the full error range.}

\FloatBarrier
\subsection{Temporal models on matched records}
The temporal benchmark asks a narrower question: whether 16 completed same-machine records improve the estimate for the next eligible event. Its history rule leaves 584 shared development targets. On this population, the current-input Ridge control achieved MAE 53.55 kWh, while TFT achieved 53.65 kWh (Table~\ref{tab:gpu}). The paired reduction relative to Ridge was $-0.105$ kWh, with a 95\% interval of [$-2.199$, 2.318] kWh. The interval includes advantages in either direction, and provides no evidence of an improvement from TFT. Adding the full history to Ridge increased MAE to 63.95 kWh; N-HiTS reached 65.28 kWh. These results concern matched event histories and should be interpreted separately from the original 800-record comparison.

\begin{table}[!htbp]\centering\small
\caption{GPU temporal benchmark on 584 identical development records. TFT and N-HiTS are three-seed prediction means; foundation-model weights remain frozen. Units are kWh per record except WAPE. These values must not be compared directly with the original 800-record ranking.}\label{tab:gpu}
\begin{tabular}{@{}lrrrr@{}}\toprule
Method & MAE & RMSE & Bias & WAPE (\%)\\\midrule
Ridge: current & 53.55 & 77.01 & $-4.72$ & 20.37\\
TFT & 53.65 & 79.11 & $-11.23$ & 20.41\\
Ridge: full window & 63.95 & 85.95 & $-1.42$ & 24.32\\
N-HiTS & 65.28 & 92.56 & $-3.05$ & 24.83\\
Chronos-2 + covariates & 85.53 & 120.82 & $-8.11$ & 32.53\\
TimesFM 2.5: energy only & 96.40 & 130.72 & $-9.22$ & 36.67\\
Chronos-2: energy only & 100.02 & 136.47 & $-11.60$ & 38.04\\
TimesFM 2.5 + XReg & 112.12 & 644.61 & $28.03$ & 42.65\\\bottomrule
\end{tabular}\end{table}


\fig{08_gpu_comparison}{gpucomparison}{1}{Accuracy and paired differences on the 584 shared development targets. (a) MAE with 95\% block-bootstrap intervals. (b) Reduction in MAE relative to current-input Ridge, so positive values favour the alternative. Pale intervals allow for seven comparisons; dark intervals are unadjusted 95\% intervals. The full interval endpoints are retained, including the wide TimesFM XReg interval. Resampling uses blocks of three observed origin dates.}

Chronos-2 benefited from the production covariates compared with its energy-only version, although its MAE of 85.53 kWh remained above the matched Ridge control. The TimesFM XReg variant had both the highest MAE and a particularly large RMSE. Three Pellet~3 records in the first validation week received predictions of approximately 2,753, 9,239 and 12,203 kWh, against observed values of 185.4, 142.4 and 70.9 kWh. In each local 16-record context, dosed mass was almost constant; the current mass lay about 1,100--1,467 historical standard deviations from the local mean. An independent float64 reconstruction closely reproduced the adapter's extreme extrapolation. All three records remain in the reported metrics.

This behavior helps explain the large gap between MAE and RMSE for that variant. It also distinguishes the local covariate adapter from the globally fitted supervised models, which estimate production effects from hundreds of records. The diagnostic concerns the tested short-context XReg configuration. It does not establish how a longer context or a separately selected, stabilized adapter would perform. The machine and week breakdowns in Figure~\ref{fig:gpustrata}, together with the complete predictions in Figure~\ref{fig:gpuall}, show the variation hidden by the overall ranking.

\fig{09_gpu_stratification}{gpustrata}{1}{MAE by machine and validation week for the same eight methods. Values are kWh per record; the labels give the number of targets in each stratum. The colour scale ends at 130 kWh to make differences among the lower errors readable, while printed values preserve every result above that limit. Machine and week panels are two views of the same 584 targets.}

\fig{10_gpu_all_predictions}{gpuall}{1}{All 4,672 predictions from the matched comparison: eight methods evaluated on 584 identical records. Colours identify the three machines and the dashed diagonal marks equality. The main axes share a 0--1,100 kWh range. The TimesFM XReg panel also includes a full-range inset containing all its predictions, including the three extremes. MAE and RMSE use the complete data in every panel.}

\FloatBarrier
\subsection{Interpreting estimates and their uncertainty}
The selected Ridge model allows each estimate to be traced to its inputs. Grouping the encoded columns by variable gives the untruncated prediction
\begin{equation}
 f_i=\widehat b+\sum_{g\in\mathcal G}C_{ig},\qquad
 C_{ig}=\sum_{k\in g}z_{ik}\widehat\beta_k,
 \label{eq:explain}
\end{equation}
where $\mathcal G$ comprises mass, batch count, hour, weekday, machine, product subfamily, presentation and bag/bulk category. This sum reproduces all 798 stored test estimates to within $1.2\times10^{-13}$ kWh; none requires non-negative clipping. Figure~\ref{fig:explanation} distinguishes the fitted coefficients from the contributions they make to actual records.

At fixed values of the other inputs, the mass coefficient corresponds to 13.30 kWh per additional tonne and the batch-count coefficient to 9.52 kWh per additional dosing batch. Their positive signs agree with an increasing energy requirement as the recorded production quantity increases. Mass and batch count are correlated, however, so the two coefficients are conditional associations rather than separately identified physical constants. The categorical effects require similar care because product, presentation and machine allocation are not randomized. In the fitted one-hot representation, the intercept of 288.09 kWh is an algebraic reference and cannot be interpreted as measured standby consumption. Small machine coefficients likewise do not establish equal intrinsic efficiency.

The real operation introduced in Figure~\ref{fig:framework} provides a concrete example. For 18.002 t dosed in five batches on Pellet~2, the model estimates 255.90 kWh and a fixed 95\% interval of [122.94, 388.85] kWh. The recorded outcome is 267.30 kWh, an excess of 11.40 kWh over the estimate and well inside the interval. This record was chosen near the median observed test energy to illustrate the calculation. Its error is not representative of every operation, as the complete residual plots show.

\fig{11_explanation}{explanation}{1}{The frozen Ridge model from coefficients to individual estimates. (a) Numerical coefficients per training standard deviation. (b) All 16 categorical coefficients in the fitted encoding. (c) Grouped contributions for every test record, with medians in red and interquartile ranges in orange; small vertical jitter separates coincident values. (d) Arithmetic decomposition of source row 2080, the example in Figure~\ref{fig:framework}. The reference combines standardized numerical means with all-zero categorical indicators and need not describe a feasible operating condition. Contributions are conditional model terms, not causal effects.}

Two exploratory extensions helped identify where modest improvements might arise (Table~\ref{tab:explore}). Allowing the mass association to vary by product subfamily reduced development MAE by 1.44 kWh, or 2.80\%. A shrunk correction from recent completed residuals reduced it by 1.71 kWh, or 3.33\%. Both retained positive intervals after adjustment for the three exploratory comparisons. By contrast, allowing mass effects to vary only by machine produced a small reduction of 0.22 kWh, with an interval that crossed zero. Figure~\ref{fig:development} shows these comparisons and their machine and week breakdowns. Product structure and stabilized residual updates are therefore candidates for a subsequent evaluation; these exploratory findings did not replace the frozen baseline.

\begin{table}[!htbp]\centering\small
\caption{Exploratory point-model extensions on the same 800 development records. Positive differences favour the extension. The adjusted intervals address these three comparisons only.}\label{tab:explore}
\begin{tabular}{@{}lrrr@{}}\toprule
Variant & MAE & MAE reduction & 98.33\% interval for reduction\\\midrule
H1: mass by machine & 51.26 & 0.217 & [$-0.082$, 0.581]\\
H2: mass by subfamily & 50.04 & 1.442 & [0.552, 2.481]\\
H3: recent residuals & 49.77 & 1.712 & [0.236, 3.156]\\\bottomrule
\end{tabular}\par\vspace{4pt}\footnotesize All quantities are in kWh per record.
\end{table}


\fig{03_development_and_hypotheses}{development}{1}{Original development results and exploratory extensions. (a) MAE of the seven original methods on 800 common records. (b) Paired MAE reductions relative to Ridge for machine-dependent mass effects (H1), product-dependent mass effects (H2) and recent-residual correction (H3), with 95\% and multiplicity-adjusted 98.33\% intervals. (c,d) Reductions by machine and validation week; positive values favour the extension. These strata describe the same development records.}

Uncertainty remained less uniform than the point comparison. The frozen 90\% and 95\% intervals covered 85.96\% and 93.23\% of test records, with mean widths of 183.63 and 258.85 kWh. Their block-bootstrap 95\% coverage intervals were [81.86, 90.24]\% and [90.53, 95.87]\%, respectively, so nominal coverage was still compatible with sampling uncertainty at the global level. Machine-level 95\% coverage was 90.00\% for Pellet~1, 96.40\% for Pellet~2 and 91.70\% for Pellet~3. Of the 54 records outside those limits, 45 were above and nine below them. These are candidates for review, without an independent diagnosis of their cause (Figure~\ref{fig:alltest}).

\fig{05_all_test_records}{alltest}{1}{Every original test record in chronological order within its machine. Black points show observed electricity, coloured points the frozen estimate and individual vertical lines the fixed 95\% intervals. Red circles identify the 54 outcomes outside their intervals. Horizontal positions are record order, so equal spacing does not imply equal elapsed time. All 798 observations and their intervals are retained, without joining adjacent records into a continuous uncertainty band.}

Updating the interval quantiles from the most recent 200 completed residuals increased coverage to 89.47\% and 94.49\%. It also increased mean width by 19.66\% and 18.85\%, while interval score improved by only 1.12\% and 0.23\% (Table~\ref{tab:intervals}). The weekly profiles in Figure~\ref{fig:uncertainty} show that coverage and width moved together. The number of records outside the 95\% interval fell from 54 to 44, but these counts cannot be interpreted as false alarms without plant adjudication. The update was explored after the original test had been examined and provides a calibration diagnostic rather than a second independent validation.

\begin{table}[!htbp]\centering\small
\caption{Fixed and post-hoc updated intervals on the same 798 archived point predictions. The updated method has not been validated on an independent period.}\label{tab:intervals}
\begin{tabular}{@{}llrrr@{}}\toprule
Method & Nominal & Coverage (\%) & Mean width (kWh) & Interval score (kWh)\\\midrule
Fixed & 90\% & 85.96 & 183.63 & 381.79\\
Updated, 200 residuals & 90\% & 89.47 & 219.74 & 377.51\\
Fixed & 95\% & 93.23 & 258.85 & 497.64\\
Updated, 200 residuals & 95\% & 94.49 & 307.64 & 496.48\\\bottomrule
\end{tabular}\end{table}


\fig{06_uncertainty_tradeoff}{uncertainty}{1}{Coverage, width and interval score for fixed calibration and the exploratory rolling update. (a,b) Machine-specific coverage against mean interval width at nominal 90\% and 95\% levels; arrows connect fixed intervals (circles) to updated intervals (triangles). (c,d) Weekly 95\% coverage and interval score, with sample sizes shown and partial weeks at the edges. Dashed horizontal lines indicate nominal coverage. Lower interval score is preferable. Both policies use the same 798 fixed point predictions.}
